% ============================================================
% CHAPITRE 7 — Sprint 4 : Entretiens, Décisions et Pilotage RH
% (fusion des anciens sprints 4 et 5 du plan initial)
% ============================================================
\chapter{Sprint 4 --- Entretiens, Décisions et Pilotage du Recrutement}

\section{Introduction}

Après le dépôt des candidatures et la présélection (Sprint~3), le
Sprint~4 couvre la \textbf{phase décisionnelle} du parcours~: le RH
retient un candidat, planifie les entretiens (en ligne puis physique),
confirme l'embauche et pilote l'activité via des tableaux de bord et
des exports.

Ce chapitre regroupe dans le rapport les contenus initialement prévus
pour les sprints~4 (\emph{Entretiens \& décisions}) et~5
(\emph{Embauche \& pilotage}). La plateforme automatise les convocations
par e-mail, centralise le calendrier RH et trace l'historique des
candidatures d'un même postulant.


\section{Objectifs du Sprint}

\begin{itemize}
  \item permettre au RH de marquer un candidat \textbf{Retenu}
        (\texttt{ACCEPTED}) et de planifier un entretien en ligne~;
  \item enregistrer l'entretien (\texttt{Interview}), horaires et lien de
        réunion issu du profil RH (Teams, Meet, etc.)~;
  \item envoyer les notifications e-mail au candidat et au RH~;
  \item planifier ensuite un \textbf{entretien physique} (convocation
        formelle avec lieu et date)~;
  \item confirmer l'\textbf{embauche} (\texttt{HIRED}) avec contrat,
        rémunération et date de prise de fonction~;
  \item afficher le \textbf{calendrier des entretiens} (\texttt{/rh/calendar})~;
  \item consulter l'\textbf{historique} des candidatures d'un candidat~;
  \item proposer un \textbf{tableau de bord} de sélection (taux de
        rétention, intégrations, répartition par agence / poste)~;
  \item exporter les candidatures (Excel mensuel et CRM complet)~;
  \item assigner un \textbf{QCM d'évaluation} aux candidats embauchés et
        générer le rapport talent (PDF).
\end{itemize}


\section{Spécification et Analyse des Besoins}

\subsection{Sprint Backlog}

\begin{table}[htbp]
  \centering
  \scriptsize
  \caption{Sprint Backlog du Sprint 4 (entretiens, embauche et pilotage)}
  \label{tab:sprint4-backlog}
  \begin{tabular}{|p{2.0cm}|c|p{4.0cm}|c|p{3.1cm}|c|}
    \hline
    \textbf{Fonctionnalité} & \textbf{ID} & \textbf{Scénario utilisateur} &
    \textbf{ID Tâche} & \textbf{Tâches} & \textbf{Est.} \\
    \hline
    Retenu en ligne
    & US-14
    & En tant que RH, je retiens un candidat et planifie un entretien en ligne.
    & T1
    & \texttt{PATCH .../status} + \texttt{InterviewService}
    & 2 j \\
    \cline{4-6}
    & & & T2 & Modal entretien + e-mails & 1,5 j \\
    \hline
    Entretien physique
    & US-15
    & En tant que RH, je planifie la réunion physique après l'entretien en ligne.
    & T3
    & Convocation physique + lieu agence
    & 1,5 j \\
    \hline
    Calendrier RH
    & US-16
    & En tant que RH, je consulte mes entretiens planifiés par date.
    & T4
    & Page \texttt{/rh/calendar}
    & 1 j \\
    \hline
    Historique candidat
    & US-17
    & En tant que RH, je consulte les candidatures passées d'un candidat.
    & T5
    & \texttt{GET /applications/history/\{id\}}
    & 1 j \\
    \hline
    Confirmation embauche
    & US-18
    & En tant que RH, je confirme l'embauche après l'entretien physique.
    & T6
    & Statut \texttt{HIRED} + e-mail confirmation
    & 2 j \\
    \hline
    Tableau de bord
    & US-19
    & En tant que RH, je visualise les indicateurs de sélection.
    & T7
    & \texttt{GET /dashboard/selection}
    & 2 j \\
    \hline
    Exports Excel
    & US-20
    & En tant que RH, j'exporte les candidatures (mensuel / CRM).
    & T8
    & \texttt{export.xlsx}, \texttt{export-full.xlsx}
    & 1,5 j \\
    \hline
    QCM embauché
    & US-21
    & En tant que RH, j'assigne un QCM post-embauche et génère le rapport.
    & T9
    & \texttt{/hired-qcm} + PDF talent
    & 2 j \\
    \hline
  \end{tabular}
\end{table}


\subsection{Diagramme de cas d'utilisation}

\begin{figure}[H]
  \centering
  \reportfig{img/uc-sprint4.jpg}
  \caption{Diagramme de cas d'utilisation du Sprint 4}
  \label{fig:uc-sprint4}
\end{figure}


\section{Description textuelle des cas d'utilisation}

\subsection{Cas d'utilisation <<~Planifier un entretien en ligne~>>}

\begin{table}[htbp]
  \centering
  \small
  \caption{Description textuelle --- Planifier un entretien en ligne}
  \label{tab:uc-online-interview}
  \begin{tabular}{|p{3.0cm}|p{10.2cm}|}
    \hline
    \textbf{Titre} & Planifier un entretien en ligne \\
    \hline
    \textbf{Acteur} & RH \\
    \hline
    \textbf{Précondition} & Candidature \texttt{SUBMITTED} ou \texttt{UNDER\_REVIEW}
          dans la zone du RH~; lien de réunion renseigné dans le profil RH (recommandé). \\
    \hline
    \textbf{Postcondition} & Statut \texttt{ACCEPTED}~; entretien \texttt{SCHEDULED}~;
          e-mails envoyés. \\
    \hline
    \textbf{Scénario nominal}
    &
    \begin{enumerate}[leftmargin=*,itemsep=0.15em,topsep=0.2em]
      \item Le RH ouvre \texttt{/rh/candidates} et clique sur \textbf{Retenu en ligne}.
      \item Il saisit la date et les horaires de début / fin.
      \item \texttt{PATCH /api/recruitment/applications/\{id\}/status} avec
            \texttt{status=ACCEPTED}, \texttt{interviewType=ONLINE}.
      \item Le service crée l'entité \texttt{Interview}, copie le lien de
            réunion du profil RH et envoie les convocations par e-mail.
      \item Le candidat voit la date et le lien dans \texttt{/jobs/applications}.
    \end{enumerate}
    \\
    \hline
    \textbf{Alternatives}
    & Lien de réunion absent~: entretien enregistré avec avertissement au RH. \\
    \hline
  \end{tabular}
\end{table}


\subsection{Cas d'utilisation <<~Planifier un entretien physique~>>}

\begin{table}[htbp]
  \centering
  \small
  \caption{Description textuelle --- Planifier un entretien physique}
  \label{tab:uc-physical-interview}
  \begin{tabular}{|p{3.0cm}|p{10.2cm}|}
    \hline
    \textbf{Titre} & Planifier un entretien physique \\
    \hline
    \textbf{Acteur} & RH \\
    \hline
    \textbf{Précondition} & Candidature déjà \texttt{ACCEPTED} avec entretien en ligne planifié. \\
    \hline
    \textbf{Postcondition} & Nouvel entretien physique planifié~; convocation formelle envoyée. \\
    \hline
    \textbf{Scénario nominal}
    &
    \begin{enumerate}[leftmargin=*,itemsep=0.15em,topsep=0.2em]
      \item Le RH clique sur \textbf{Réunion physique} depuis la fiche candidat.
      \item Il choisit date, horaires et éventuellement le lieu (sinon adresse agence).
      \item Le backend annule l'entretien précédent et crée un entretien
            \texttt{PHYSICAL}.
      \item Un e-mail de convocation (modèle DAAM) est envoyé au candidat et au RH.
    \end{enumerate}
    \\
    \hline
    \textbf{Alternatives}
    & Tentative sans entretien en ligne préalable~: refus métier. \\
    \hline
  \end{tabular}
\end{table}


\subsection{Cas d'utilisation <<~Confirmer une embauche~>>}

\begin{table}[htbp]
  \centering
  \small
  \caption{Description textuelle --- Confirmer une embauche}
  \label{tab:uc-hire}
  \begin{tabular}{|p{3.0cm}|p{10.2cm}|}
    \hline
    \textbf{Titre} & Confirmer une embauche \\
    \hline
    \textbf{Acteur} & RH \\
    \hline
    \textbf{Précondition} & Candidature \texttt{ACCEPTED} avec entretien physique planifié. \\
    \hline
    \textbf{Postcondition} & Statut \texttt{HIRED}~; données contrat enregistrées~;
          e-mail de confirmation envoyé. \\
    \hline
    \textbf{Scénario nominal}
    &
    \begin{enumerate}[leftmargin=*,itemsep=0.15em,topsep=0.2em]
      \item Le RH clique sur \textbf{Embauche} depuis \texttt{/rh/candidates}.
      \item Il renseigne date de prise de fonction, salaire net, type de contrat,
            horaires et avantages.
      \item \texttt{PATCH .../status} avec \texttt{status=HIRED} et champs d'embauche.
      \item Le candidat apparaît dans la liste des embauchés et le tableau de bord
            est mis à jour.
    \end{enumerate}
    \\
    \hline
  \end{tabular}
\end{table}


\subsection{Cas d'utilisation <<~Consulter le tableau de bord~>>}

\begin{table}[htbp]
  \centering
  \small
  \caption{Description textuelle --- Consulter le tableau de bord}
  \label{tab:uc-dashboard}
  \begin{tabular}{|p{3.0cm}|p{10.2cm}|}
    \hline
    \textbf{Titre} & Consulter le tableau de bord de sélection \\
    \hline
    \textbf{Acteur} & RH \\
    \hline
    \textbf{Précondition} & RH authentifié avec zones affectées. \\
    \hline
    \textbf{Postcondition} & Indicateurs affichés (candidats, retenus, intégrés, taux). \\
    \hline
    \textbf{Scénario nominal}
    &
    \begin{enumerate}[leftmargin=*,itemsep=0.15em,topsep=0.2em]
      \item Le RH ouvre \texttt{/rh/dashboard}.
      \item Il filtre éventuellement par mois / année.
      \item \texttt{GET /api/recruitment/dashboard/selection} retourne les statistiques
            agrégées par zone, agence, poste et responsable.
      \item Les graphiques et cartes KPI sont affichés.
    \end{enumerate}
    \\
    \hline
  \end{tabular}
\end{table}


\subsection{Autres cas d'utilisation}

\begin{itemize}
  \item \textbf{Consulter le calendrier des entretiens}~: vue mensuelle avec
        liste des entretiens du jour (\texttt{/rh/calendar}).
  \item \textbf{Consulter l'historique candidat}~: toutes les candidatures
        passées d'un même utilisateur.
  \item \textbf{Suivre sa convocation} (candidat)~: date, horaire et lien de
        réunion dans \texttt{/jobs/applications}.
  \item \textbf{Exporter les candidatures}~: Excel mensuel par feuille de mois
        ou export CRM complet.
  \item \textbf{Assigner un QCM embauché}~: évaluation post-recrutement et
        génération du rapport talent PDF.
  \item \textbf{Annuler un entretien}~: lors d'un refus (\texttt{REJECTED}) ou
        d'un désistement (\texttt{DESISTED}), l'entretien planifié est annulé.
\end{itemize}


\section{Conception}

\subsection{Diagramme de classes}

Le domaine du Sprint~4 enrichit \textbf{JobApplication} avec les champs
d'entretien et d'embauche, et introduit l'entité \textbf{Interview} ainsi
que le module \textbf{HiredQcmAssignment} pour l'évaluation post-embauche.

\begin{figure}[H]
  \centering
  \reportfig{img/class-sprint4.jpg}
  \caption{Diagramme de classes du Sprint 4}
  \label{fig:class-sprint4}
\end{figure}


\subsection{Diagrammes de séquence}

\subsubsection{Diagramme de séquence <<~Planifier un entretien en ligne~>>}

\begin{figure}[H]
  \centering
  \reportfig{img/seq-online-interview.jpg}
  \caption{Diagramme de séquence --- Planifier un entretien en ligne}
  \label{fig:seq-online-interview}
\end{figure}

\subsubsection{Diagramme de séquence <<~Confirmer une embauche~>>}

\begin{figure}[H]
  \centering
  \reportfig{img/seq-hire.jpg}
  \caption{Diagramme de séquence --- Confirmer une embauche}
  \label{fig:seq-hire}
\end{figure}

\subsubsection{Diagramme de séquence <<~Consulter le tableau de bord~>>}

\begin{figure}[H]
  \centering
  \reportfig{img/seq-dashboard.jpg}
  \caption{Diagramme de séquence --- Consulter le tableau de bord}
  \label{fig:seq-dashboard}
\end{figure}


\section{Réalisation}

\subsection{Interface RH --- décisions et entretiens}

La page \texttt{/rh/candidates} propose le flux décisionnel~: \textbf{Retenu
en ligne}, \textbf{Réunion physique}, \textbf{Embauche}, ainsi que
\textbf{Non retenu} et \textbf{Désisté}. Un modal permet de saisir la date
et les horaires avant envoi au backend.

\begin{figure}[H]
  \centering
  \reportfig{img/ui-rh-decisions.png}
  \caption{Interface de traitement des candidatures --- décisions RH}
  \label{fig:ui-rh-decisions}
\end{figure}

\begin{figure}[H]
  \centering
  \reportfig{img/ui-interview-modal.png}
  \caption{Interface de planification d'un entretien}
  \label{fig:ui-interview-modal}
\end{figure}


\subsection{Calendrier et historique}

Le calendrier RH regroupe les entretiens planifiés. L'historique candidat
permet au RH de consulter les candidatures antérieures avant de statuer.

\begin{figure}[H]
  \centering
  \reportfig{img/ui-calendar.png}
  \caption{Interface du calendrier des entretiens (RH)}
  \label{fig:ui-calendar}
\end{figure}


\subsection{Confirmation d'embauche}

Après l'entretien physique, le RH confirme l'embauche via un formulaire
(contractuel) qui déclenche l'e-mail de confirmation au candidat.

\begin{figure}[H]
  \centering
  \reportfig{img/ui-hire-modal.png}
  \caption{Interface de confirmation d'embauche}
  \label{fig:ui-hire-modal}
\end{figure}


\subsection{Tableau de bord et exports}

Le tableau de bord \texttt{/rh/dashboard} synthétise les campagnes, le taux
de sélection, les intégrations du mois et les répartitions par agence et
par poste. Les exports Excel complètent le pilotage opérationnel.

\begin{figure}[H]
  \centering
  \reportfig{img/ui-dashboard.png}
  \caption{Interface du tableau de bord de sélection (RH)}
  \label{fig:ui-dashboard}
\end{figure}


\subsection{Espace candidat --- convocation}

Le candidat retenu consulte la date de son entretien et le lien de réunion
depuis \texttt{/jobs/applications}.

\begin{figure}[H]
  \centering
  \reportfig{img/ui-candidate-interview.png}
  \caption{Interface candidat --- convocation à l'entretien}
  \label{fig:ui-candidate-interview}
\end{figure}


\section{Conclusion}

Le Sprint~4 a complété le parcours métier après la candidature~: décision
RH, entretiens en ligne et physique, confirmation d'embauche, pilotage par
tableau de bord et exports. Le workflow impose une progression logique
(en ligne $\rightarrow$ physique $\rightarrow$ embauche) et automatise les
communications par e-mail.

Le chapitre suivant présentera le \textbf{module Réclamations et
Notifications}, microservice complémentaire prévu dans l'architecture
globale de la plateforme.

\clearpage
